\section{Specialization beyond a first-order weight}
\label{rev:scarcity-refinements}

The main separation theorem exploits cancellation of the common Tate
period at two parameters. A complementary argument works at one
parameter by taking a field norm of the period relation. Its value is
the uniformity in arbitrary polynomial weights: their derivatives
remain in one symmetric-square differential module.


We use the following fixed-field relation theorem from the preprint
\cite[Theorem~4.2(3)]{BaldiBinyaminiUrbanik2026}. Its uniformity concerns
the degree of a relation, not the heights of its coefficients.

\begin{theorem}[Baldi--Binyamini--Urbanik, Theorem 4.2(3)]
\label{thm:C-BBU}
Let $L$ be a rank-$n$ $G$-operator over a number field $K$ on
$\mathbb P^1$, with singular set $\Sigma$.  Choose any complex
fundamental matrix $\mathcal P$ of its solutions and a branch near
$x=\infty$.  Let $B_{\mathcal P}$ be the $K$-Zariski closure of its
graph.  A polynomial relation at $x\in\mathbb P^1(K)$ is
\emph{nontrivial} if it does not vanish identically on any
$K$-irreducible component of the fiber $B_{\mathcal P,x}$ through
$\mathcal P(x)$.

For every fixed integer $\delta\geq1$ and every integer $\nu\geq1$
there exist an integer
\[
 1\leq D\leq n^2(|\Sigma|+1)^2+\nu
\]
and a constant $\kappa$ such that, for every integer $t>\kappa$,
at least $\nu$ of the $n^2+\nu$ points
\[
 t,t+D,\ldots,t+(n^2+\nu-1)D
\]
satisfy no nontrivial polynomial relation over $K$ of total degree
at most $\delta$ among the entries of $\mathcal P$.
\end{theorem}

The theorem puts no bound on the heights of the relation coefficients.
The fundamental matrix need not have algebraic initial values.  The
fixed field $K$ is essential; below a norm polynomial brings every
multiplier of bounded degree into the single field $K=\mathbb Q$.

\begin{theorem}[Uniformly sparse denominators]
\label{thm:C-density-zero}
Fix a standard level $\ell\in\{1,2,3,4\}$ and an integer $d\geq1$.
Let $\mathcal E_{\ell,d}$ consist of the nonzero integers $C$ for
which there exist integers $A,B$, not both zero, and an algebraic
number $\alpha$ such that the standard series converges ordinarily,
\[
 (A+B\thetaop)F_{s_\ell}(R_\ell/C)=\frac{\alpha}{\pi},
 \qquad [\mathbb Q(\alpha^2):\mathbb Q]\leq d.
\]
Then $\mathcal E_{\ell,d}$ has upper Banach density zero:
\[
 \limsup_{N\to\infty}\ \sup_{M\in\mathbb Z}
 \frac{\#(\mathcal E_{\ell,d}\cap\{M+1,\ldots,M+N\})}{N}=0.
\]
In particular,
\[
 \#\{C\in\mathcal E_{\ell,d}:0<|C|\leq X\}=o(X).
\]
More precisely, for each sign $\epsilon\in\{1,-1\}$ and each
$\nu\geq1$ there exist $D\leq64+\nu$ and $\kappa$ such that every
integer $t>\kappa$ satisfies
\[
 \#\{0\leq j\leq\nu+3:
       \epsilon(t+jD)\in\mathcal E_{\ell,d}\}\leq4.
\]
Neither $D$ nor $\kappa$ depends on $A,B$, the field generated by
$\alpha$, or the height of $\alpha$.  They may depend on
$\ell,d,\nu,\epsilon$ and the chosen branch.  In particular, the
density-zero conclusion holds under the more familiar restriction
$[\mathbb Q(\alpha):\mathbb Q]\leq d$.
\end{theorem}

\begin{proof}
Put $s=s_\ell$, $R=R_\ell$, $a=s/2$, and $b=(1-s)/2$.  Thus
$a+b=1/2$, and Clausen's formula gives
\[
 f(z)={}_2F_1(a,b;1;z),\qquad F_s(z)=f(z)^2.
\]
The function $f$ is annihilated by
\begin{equation}\label{rev:eq:gauss-operator-appendix}
 L_s=z(1-z)\partial_z^2+
          (1-\tfrac32z)\partial_z-ab.
\end{equation}
It is a rank-two $G$-operator: the rational hypergeometric germ is a
$G$-function, its minimal equation has rank two by
Lemma~\ref{rev:lem:product-monodromy}, and factors of $G$-operators
remain $G$-operators \cite[Theorem~3.5]{BaldiBinyaminiUrbanik2026}.

Choose a fundamental matrix
\[
 \mathcal P(z)=
 \begin{pmatrix}p_1(z)&p_2(z)\\p'_1(z)&p'_2(z)\end{pmatrix},
 \qquad p_1(z)=\pi f(z),
\]
and scale the second solution so that
\begin{equation}\label{eq:C-wronskian}
 \det\mathcal P(z)=\frac{\pi}{z\sqrt{1-z}}.
\end{equation}
This is possible because Abel's identity for
\eqref{rev:eq:gauss-operator-appendix} says that every nonzero Wronskian is a
constant multiple of $1/(z\sqrt{1-z})$.  Here $\sqrt{1-z}$ is the
branch equal to one at zero.  The same slit punctured neighborhood
of zero may contain both real rays; alternatively the two signs
may be treated with separate branch choices.

We first prove that the $\overline{\mathbb Q}$-Zariski closure of
the graph of $\mathcal P$ is the whole
$\mathbb P^1\times\operatorname{GL}_2$.  The connected monodromy is $\operatorname{SL}_2$ by
Lemma~\ref{rev:lem:product-monodromy}.

For completeness, this last assertion about monodromy does not by
itself give the claimed arithmetic graph closure: its determinant
must also be treated.  Suppose a polynomial over
$\overline{\mathbb Q}(z)$ vanishes on $\mathcal P(z)$.  Analytic
continuation makes it vanish on every
$\operatorname{SL}_2$-orbit of $\mathcal P(z)$, hence on the
hypersurface with determinant $\det\mathcal P(z)$.  Localize the
matrix coordinate ring at $X_{11}$ and replace $X_{22}$ with
$(T+X_{12}X_{21})/X_{11}$, where $T=\det X$.  Expanding in
$X_{11}^{\pm1},X_{12},X_{21}$ shows that each resulting coefficient
polynomial in $T$ vanishes at
$T=\pi/(z\sqrt{1-z})$.  This value is transcendental over
$\overline{\mathbb Q}(z)$: otherwise the constant $\pi$ would be
algebraic over that field, hence algebraic over
$\overline{\mathbb Q}$, a contradiction.  All coefficient
polynomials vanish identically.  There was therefore no nonzero
polynomial relation on the graph.  In particular its
$\mathbb Q$-Zariski closure has the same property, and every
nonzero polynomial is nontrivial on its smooth fibers.

Let $z=R/C$ with $z\neq0,1$.  Writing a matrix variable as
$X=(X_{ij})$, define homogeneous degree-two and degree-four
polynomials
\[
 U_C(X)=A X_{11}^2+2BzX_{11}X_{21},\qquad
 V_C(X)=z^2(1-z)(\det X)^2.
\]
The assumed series identity and \eqref{eq:C-wronskian} give
\begin{equation}\label{eq:C-norm-relation}
 U_C(\mathcal P(z))=\alpha\pi,
 \qquad V_C(\mathcal P(z))=\pi^2,
 \qquad U_C(\mathcal P(z))^2=\alpha^2 V_C(\mathcal P(z)).
\end{equation}
Let $L=\Q(\alpha^2)$, with $e=[L:\Q]\le d$. Descent to the fixed
coefficient field is the norm
\begin{equation}\label{eq:C-rational-norm}
 Q_C(X)=\operatorname{Nm}_{L/\Q}\bigl(U_C(X)^2-\alpha^2V_C(X)\bigr).
\end{equation}
This polynomial
has rational coefficients, total degree $4e\leq4d$, and vanishes
at $\mathcal P(z)$.  It is nonzero.  Indeed, $U_C$ is nonzero
because $(A,B)\neq(0,0)$ and $z\neq0$.  It cannot be proportional
to $\det X$, whose $X_{12}X_{21}$ term is nonzero.  Consequently
$U_C^2/V_C$ is a nonconstant rational function on
$\operatorname{GL}_2$, and no nonzero polynomial $m$ with
constant coefficients annihilates it.  Thus
\eqref{eq:C-rational-norm} is a nontrivial relation in the exact
sense of Theorem~\ref{thm:C-BBU}.

Use the rational base coordinate $x=C=R/z$.  The pullback of
$L_s$ remains a $G$-operator by
\cite[Theorem~3.5(4)]{BaldiBinyaminiUrbanik2026}; it has rank two
and singular set $\{0,R,\infty\}$, of cardinality three.  One may
retain the displayed matrix by choosing the differential-module
basis corresponding to $1,\partial_z$; equivalently, changing to
$1,\partial_x$ is an invertible rational row change and does not
affect the graph-closure or polynomial-degree argument.  Apply
Theorem~\ref{thm:C-BBU} with $K=\mathbb Q$, $\delta=4d$, and
$n=2$.  This gives the asserted positive-denominator block bound.
Replacing the base coordinate by $x=-C$ gives the negative bound,
with singular set $\{0,-R,\infty\}$.  The excluded points
$C=0,R$ are finite and have no effect on the conclusion.  The
conditionally convergent point $z=-1$ is nonsingular and is covered
by analytic continuation and Abel's theorem, or simply omitted as
one additional finite exception.

Finally fix $\nu$ and the resulting $D$.  In any interval of $N$
consecutive positive integers, discard its intersection with
$[1,\kappa]$ and partition each residue class modulo $D$ into
blocks of $\nu+4$ terms, leaving fewer than $\nu+4$ terms per
residue class.  Every full block contains at most four exceptional
terms.  Thus the number of exceptional terms is at most
\[
 \frac{4N}{\nu+4}+D(\nu+4)+\max(0,\lceil\kappa\rceil).
\]
The error is independent of the position of the interval.  Apply
the same argument to negative denominators.  Splitting any integer
interval at zero therefore gives an upper Banach density bound
$4/(\nu+4)$.  Letting $\nu$ increase proves the claim.
\end{proof}

\begin{proposition}[A degree bound for two nearby identities]
\label{prop:C-paired-degree-growth}
Fix two standard levels, a common sign, and a gap $D\geq0$,
excluding the identical-level, zero-gap case.  There is a constant
$K>0$, depending only on these fixed data, such that any two
nonzero identities at denominators $\epsilon t$ and
$\epsilon(t+D)$, with $t$ sufficiently large and respective
multiplier-square degrees
\[
 d_i=[\mathbb Q(\alpha_i^2):\mathbb Q],
\]
must satisfy
\[
 \log t<K(d_1d_2)^4\log(d_1d_2+1).
\]
In particular, for a suitable $c>0$,
\[
 d_1d_2\geq c\left(\frac{\log t}{\log\log t}\right)^{1/4},
 \qquad
 \max(d_1,d_2)\geq c
       \left(\frac{\log t}{\log\log t}\right)^{1/8}.
\]
One may replace $d_1d_2$ in the first inequality by
$[\mathbb Q(\alpha_1^2/\alpha_2^2):\mathbb Q]$, adjusting $K$.
\end{proposition}

\begin{proof}
For $Y=F_s(z)$ and $Z=\thetaop F_s(z)$, the rank-two Gauss equation
gives the rational differential relations
\[
 \thetaop Y=Z,\qquad
 \thetaop Z=\frac{Z^2}{2Y}
             +\frac{zZ}{2(1-z)}
             +\frac{s(1-s)zY}{2(1-z)}.
\]
Rational composition preserves rational differential closure.
Consequently the differential field generated by the four
$G$-functions in the cross-level proof has transcendence degree
$\mu'=4$ over $\mathbb Q(u)$; their algebraic independence gives
equality rather than just an upper bound.

With this $\mu'$, the height condition in
\cite[Theorem~3.13, (3.5)]{BaldiBinyaminiUrbanik2026} is
$h\geq c_1\delta^4(\log\delta+1)$.  Its smallness condition at
$u=1/t$, where $h=\log t$, is implied by
\[
 h^{4/5}\geq c_2'\delta^{4/5}\log h.
\]
Both conditions therefore hold whenever
$h\geq K_0\delta^4\log(\delta+1)$, with a fixed sufficiently
large $K_0$.  For example, after increasing a fixed lower bound
on $h$, the inequality $\log h\leq h^{3/10}$ shows that the
second condition follows already from
$h\geq K_1\delta^{8/5}$; the stated bound dominates this for
every $\delta\geq1$.

The polynomial constructed from the multiplier-square ratio has
degree at most $2d_1d_2$.  If the displayed sufficient height bound
held for $\delta=2d_1d_2$, no such polynomial could vanish.
Its contrapositive, with constants absorbed into $K$, proves the
first assertion.  The two lower bounds follow by elementary
inversion; if $d_1d_2>\log t$ they are immediate, while otherwise
$\log(d_1d_2+1)\leq\log(\log t+1)$.  Using the actual degree of
the ratio's minimal polynomial gives the final strengthening.
\end{proof}

\paragraph{What this does and does not prove.}
All 36 listed standard identities have $\alpha^2\in\mathbb Q$,
so the theorem already excludes identities of that multiplier type
at a density-one set of denominators, with no bound on $A$ or $B$.
The fixed-gap proposition further shows that such identities
cannot continue to occur in clusters of bounded diameter at
arbitrarily large denominators.  These results do not prove that
the exceptional denominators are finite, identify a particular
denominator as nonexceptional, or force CM at an individual
exceptional denominator.
The union over all $d$ cannot be taken inside the density
conclusion: a countable union of density-zero sets can have density
one.  In particular the original unrestricted algebraic multiplier
problem remains open in this argument.  The cited input is a
recent preprint; this deduction is unconditional relative to its
stated theorem, without a transcendence conjecture.

The density conclusion also holds for the union over the four
standard levels, since this is a finite union.  The separate
cross-level proof establishes the stronger divergent-gap property
for that union.  Consequently, if a set of integer
denominators of positive upper Banach density supports standard
identities with selected algebraic multipliers $\alpha_C$, then
the degrees $[\mathbb Q(\alpha_C^2):\mathbb Q]$ cannot be
uniformly bounded.  More precisely, for every fixed $d$, the
denominators for which some identity has degree at most $d$ form
a set of upper Banach density zero.

\subsection{All polynomial weights and fixed elliptic pullbacks}
\label{sec:C-polynomial-density}

The rank-two formulation gives more than a density theorem for linear
weights. It controls every rational polynomial weight at once, without
bounding its degree or coefficient height. The multiplier must be
nonzero; otherwise Picard--Fuchs annihilators give zero identities at
every parameter.

\begin{lemma}[Quadratic-form exclusion]\label{lem:C-quadratic-exclusion}
Let $L$ be a fixed rank-two $G$-operator over $\Q(t)$, with fundamental
matrix $X(t)$, such that its arithmetic graph closure is
$\mathbb P^1\times\operatorname{GL}_2$ and
\[
 (\det X(t))^2=\pi^2r(t),\qquad r\in\Q(t)^\times.
\]
Fix positive integers $e,d$. Consider regular reciprocal integer
parameters $t=1/N$ at which
\begin{equation}\label{eq:C-quadratic-specialization}
 Q_N(X_{11}(t),X_{21}(t))=\alpha\pi
\end{equation}
for a homogeneous quadratic form $Q_N$ whose coefficients belong to
some number field of degree at most $e$, and some
$\alpha\in\Qbar^\times$ with $[\Q(\alpha^2):\Q]\le d$.
These positive integers $N$ have natural density zero. For every
$\nu\ge1$ there are integers
\[
 1\le D\le4(|\Sigma(L)|+1)^2+\nu
\]
and a constant $\kappa$ such that every block
$N,N+D,\ldots,N+(\nu+3)D$, with $N>\kappa$, contains at least
$\nu$ parameters admitting no identity
\eqref{eq:C-quadratic-specialization}. The analogous assertion holds
for negative integers, ordered by absolute value.
\end{lemma}
\begin{proof}
Let $E_N$ contain the coefficients of $Q_N$ and put
$T_N=E_N\Q(\alpha^2)$, so $[T_N:\Q]\le ed$. Square
\eqref{eq:C-quadratic-specialization} and take the coefficient-field
norm of
\[
 r(1/N)Q_N(X_{11},X_{21})^2-\alpha^2(\det X)^2.
\]
This gives a rational polynomial relation of degree at most $4ed$.
Each conjugate factor is nonzero on $\operatorname{GL}_2$: its
quadratic-form term depends only on the first column, whereas its
nonzero determinant term varies with the second column. The product
is therefore nonzero and, by the graph-closure hypothesis, nontrivial
in the sense of Theorem~\ref{thm:C-BBU}. That theorem gives the block
assertion after the reciprocal change of coordinate.

Write $k=\nu+4$. Summing the bound of at most four exceptional entries
per block over all starting integers up to $Y$ counts each exceptional
integer $k$ times, apart from $O(kD)$ boundary terms. The upper density
is thus at most $4/k$. Letting $\nu$ increase proves density zero.
The same count is uniform on arbitrary intervals, so the upper Banach density is also zero. Replacing $t$ by $-t$ proves the negative version.
\end{proof}

If $\det X=\pi w(t)$ with $w\in\Q(t)^\times$, one may instead norm
$w(1/N)Q_N-\alpha\det X$. A bound $[\Q(\alpha):\Q]\le d$ then
requires relation degree only $2ed$.

\begin{theorem}[Unrestricted polynomial weights]
\label{thm:C-all-polynomial-weights}
Under the matrix hypotheses of Lemma~\ref{lem:C-quadratic-exclusion},
suppose an analytic germ has the representation
\[
 H(t)=k(t)X_{11}(t)^2/\pi^2,
\]
where $k$ lies in a fixed algebraic extension of $\Q(t)$ of degree
at most $e$. Assume $X'=M(t)X$ with $M\in M_2(\Q(t))$.
For each fixed $d$, the reciprocal integer parameters admitting
\[
 P\left(t\frac d{dt}\right)H(1/N)=\frac\alpha\pi,
 \qquad P\in\Q[T],\quad
 0\ne\alpha\in\Qbar,\quad [\Q(\alpha^2):\Q]\le d,
\]
have density zero and satisfy the block conclusion of
Lemma~\ref{lem:C-quadratic-exclusion}. The polynomial $P$ may vary
without any restriction on its degree or coefficient height.
At points of ordinary interior convergence this includes all
polynomially weighted series
$\sum_{n\ge0}h_nP(n)N^{-n}$ for $H(t)=\sum h_nt^n$.
\end{theorem}
\begin{proof}
Differentiation preserves the three-dimensional space of homogeneous
quadratic forms in the first column of $X$. Its coefficients remain
in the fixed field $\Q(t,k)$, since differentiation of an algebraic
function stays in the same function field. Induction gives
\[
 \pi^2\left(t\frac d{dt}\right)^jH
 =Q_j(t;X_{11},X_{21}),\qquad j\ge0,
\]
with $Q_j$ homogeneous quadratic. Outside a fixed finite pole and
ramification locus, all coefficients specialize into one field of
degree at most $e$. Their linear combination for any $P$ therefore
satisfies Lemma~\ref{lem:C-quadratic-exclusion}. The argument does not
bound $j$, and the excluded locus is independent of $j$.
\end{proof}

For the four standard families this assertion can be seen directly.
Put $f={}_2F_1(s/2,(1-s)/2;1;z)$, $v=\thetaop f$, and
\[
 c=\frac{s(1-s)z}{4(1-z)},\qquad b_0=\frac z{2(1-z)}.
\]
The Gauss equation gives $\thetaop v=cf+b_0v$, hence
\[
 \thetaop(f^2)=2fv,\qquad
 \thetaop(fv)=cf^2+b_0fv+v^2,\qquad
 \thetaop(v^2)=2cfv+2b_0v^2.
\]
Every polynomial weight reduces to a rational quadratic form in
$(f,v)$. Thus Theorem~\ref{thm:C-density-zero} holds simultaneously
for all rational polynomial weights with a nonzero multiplier of
bounded square-degree.

\begin{remark}[A necessary nonvanishing hypothesis]
For any rational $z_0$ with $0<|z_0|<1$, the nonzero polynomial
\[
 P_{z_0}(T)=T^3-z_0(T+\tfrac12)(T+s)(T+1-s)
\]
has an absolutely convergent weighted series of value zero. This is
the hypergeometric differential equation
\[
 [\thetaop^3-z(\thetaop+\tfrac12)(\thetaop+s)
                   (\thetaop+1-s)]F_s=0
\]
specialized at $z_0$. Clearing denominators and dividing their gcd
gives a primitive integral polynomial weight. Allowing $\alpha=0$
in Theorem~\ref{thm:C-all-polynomial-weights} would therefore make
every rational parameter exceptional. This does not affect the
original first-order conjecture.
\end{remark}

\begin{corollary}[Gauges and rational pullbacks]
\label{cor:C-pullback-density}
Let $\phi\in\Q(t)$ be nonconstant of degree $r$ and let $g$ be a
fixed algebraic function of degree at most $e$ over $\Q(t)$. For a
specified analytic germ
\[
 H(t)=g(t)F_s(\phi(t)),
 \qquad s\in\{\tfrac16,\tfrac14,\tfrac13,\tfrac12\},
\]
Theorem~\ref{thm:C-all-polynomial-weights} holds, with relation-degree
bound $4ed$ and block step
\[
 D\le4(3r+1)^2+\nu.
\]
One omits the finite set of undefined matrix specializations.
In particular rational Euler gauges leave the underlying rank and
singularity bound unchanged.
\end{corollary}
\begin{proof}
A rational pullback is a $G$-operator by
\cite[Theorem~3.5(4)]{BaldiBinyaminiUrbanik2026}. After deleting branch
values it induces a finite-index subgroup on fundamental groups, so
the connected algebraic monodromy remains $\operatorname{SL}_2$.
The determinant in the derivative frame is the old determinant
composed with $\phi$, multiplied by $\phi'$. Its square is still
$\pi^2$ times a nonzero rational function. The monodromy and
transcendental-determinant argument in the proof of
Theorem~\ref{thm:C-density-zero} therefore gives the full arithmetic
graph closure.

The non-apparent singularities lie above $\{0,1,\infty\}$ and number
at most $3r$. A critical point above an ordinary point has two
holomorphic pulled-back solutions and is apparent in the scalar
frame, so does not enlarge the set of Definition~3.7 of the cited
source. The algebraic gauge changes the quadratic coefficients, not
the rank-two differential module. Apply
Theorem~\ref{thm:C-all-polynomial-weights}.
\end{proof}

The safe bounds for the Euler, quadratic-companion, Domb, and
level-nine transformations are respectively
$64+\nu$, $196+\nu$, $400+\nu$, and $1444+\nu$.
The level-nine gauge has degree two, giving relation-degree bound
$8d$; its rational parameter map has degree six.
These are densities of integer parameters. They make no claim about
the prime density of raw weighted Euler truncations at one fixed
parameter.

\subsection{A general elliptic-family version}
\label{sec:C-elliptic-density}

\begin{proposition}\label{prop:C-general-elliptic-density}
Let $E\to U$ be a fixed nonisotrivial elliptic family over $\Q$, where
$U$ is a nonempty open subset of $\mathbb P^1$. Choose an algebraic
de Rham frame over $\Q(t)$ and one locally flat nonzero integral
cycle, whose holomorphic period is $\omega(t)$. Suppose
\[
 H(t)=k(t)\omega(t)^2/\pi^2,
\]
with $k$ algebraic of degree at most $e$ over $\Q(t)$.
For each $d$, the integers $N$ admitting a rational polynomial weight
and a nonzero multiplier $\alpha$ of degree at most $d$ with
$P(t\,d/dt)H(1/N)=\alpha/\pi$ have density zero. A degree-$2ed$
relation and a fixed rank-two geometric $G$-operator suffice.
\end{proposition}
\begin{proof}
Extend the chosen nonzero integral cycle to a rational Betti basis whose intersection pairing is one. This is possible even for a nonprimitive first cycle, by dividing the second cycle by the first cycle's divisibility. Normalize the cup product of the de Rham basis so that its period
matrix has determinant $2\pi i$, and multiply its second Betti column
by $1/(2i)$. The resulting fundamental matrix has determinant $\pi$,
while its first column is unchanged. It is a matrix for the same
rank-two differential module; arbitrary complex solution bases are
permitted in the cited theorem. A cyclic scalar equation for this
Gauss--Manin module is a $G$-operator by
\cite[Theorem~3.10]{BaldiBinyaminiUrbanik2026}.

Nonisotriviality gives connected monodromy $\operatorname{SL}_2$:
after a finite level-structure cover, the nonconstant moduli map is
a finite map of curves over suitable open subsets, giving
finite-index modular monodromy. The graph-closure proof used above
now applies with the constant determinant $\pi$. More explicitly,
a rational polynomial relation on the graph vanishes on each full
$\operatorname{SL}_2$ orbit. At an algebraic regular base point it
therefore vanishes on the determinant-$\pi$ hypersurface; the
transcendence of $\pi$ forces that specialization to be the zero
polynomial. Infinitely many such base points force the original
polynomial to vanish identically. The graph closure is the full
product. Use the degree-two norm variant of
Lemma~\ref{lem:C-quadratic-exclusion} and the same quadratic-form
reduction.
\end{proof}

This proposition applies to any higher-level elliptic period-square
realization on a rational base after its normalization is verified.
It supplies no realization for general triangle-group or compact
Shimura functions. Nor does density zero imply finiteness, identify
the exceptional parameters, or control the union over all multiplier
degrees. The individual non-CM exclusion remains the unresolved step.




\begin{corollary}[Divergent gaps for all polynomial weights]
\label{cor:C-polynomial-divergent-gaps}
For $d\ge1$, let $\mathcal E^{\mathrm{poly}}_{\ell,d}$ be the integer
denominators $C$ in the interior convergence domain at which
\[
 P(\thetaop)F_{s_\ell}(R_\ell/C)=\frac\alpha\pi,
 \qquad P\in\Q[T],\qquad
 0\ne\alpha\in\Qbar,\qquad [\Q(\alpha^2):\Q]\le d.
\]
For each fixed positive integer $D$, there are only finitely many
positive $t$ for which both $\epsilon t$ and $\epsilon(t+D)$ lie
in $\mathcal E^{\mathrm{poly}}_{\ell,d}$, for either fixed sign
$\epsilon$. The same conclusion holds with two different standard
levels, also allowing $D=0$. Consequently the successive gaps in the
positive part, and in the absolute values of the negative part, of
$\bigcup_{\ell=1}^4\mathcal E^{\mathrm{poly}}_{\ell,d}$ tend to
infinity. Neither polynomial degree nor coefficient height is bounded.
\end{corollary}
\begin{proof}
Fix the levels, sign, and gap, excluding the equal-level zero-gap
case. Write $r_i=\epsilon R_{\ell_i}$ and
$f_i(z)={}_2F_1(s_{\ell_i}/2,(1-s_{\ell_i})/2;1;z)$.
The four genuine $G$-functions
\[
 f_1(r_1u),\quad (\thetaop f_1)(r_1u),\quad
 f_2\!\left(\frac{r_2u}{1+Du}\right),\quad
 (\thetaop f_2)\!\left(\frac{r_2u}{1+Du}\right)
\]
are holomorphic at zero and algebraically independent over $\Q(u)$.
Indeed the product-monodromy proof of
Theorem~\ref{thm:C-absolute-gaps} acts independently on their
two nonzero rank-two solution columns, and each
$\operatorname{SL}_2$ orbit is dense in its affine two-space.
Here $\thetaop$ differentiates the displayed hypergeometric argument.

At $u=1/t$, any two polynomially weighted identities reduce by the
quadratic recursions above to
\[
 Q_1(G_1,G_2)=\alpha_1/\pi,\qquad
 Q_2(G_3,G_4)=\alpha_2/\pi,
\]
with rational homogeneous quadratic forms $Q_i$. They are nonzero
because $\alpha_i\ne0$. Put
$\beta=\alpha_1^2/\alpha_2^2$, so $[\Q(\beta):\Q]\le d^2$.
Writing $m_\beta(T)=\sum_{j=0}^e c_jT^j$ for its minimal polynomial,
the polynomial
\[
 \sum_{j=0}^e c_jQ_1(T_1,T_2)^{2j}
                    Q_2(T_3,T_4)^{2(e-j)}
\]
has rational coefficients, total degree at most $4d^2$, and vanishes
at the four specialized $G$-functions. It is nonzero: each factor
$Q_1^2-\sigma(\beta)Q_2^2$ involves nonzero homogeneous forms in
disjoint variable sets and $\sigma(\beta)\ne0$.
Lemma~\ref{lem:C-reciprocal-specialization} excludes this relation
for all sufficiently large $t$, uniformly in its coefficients.
Taking the maximum of the finitely many thresholds for ordered level
pairs and $0\le D\le H$ gives the asserted divergent-gap statement.
\end{proof}

